\documentclass[11pt,a4paper]{article}
\usepackage{fontspec,amsmath,amssymb,amsthm,microtype}
\setmainfont{lmroman10-regular.otf}[BoldFont=lmroman10-bold.otf,ItalicFont=lmroman10-italic.otf,BoldItalicFont=lmroman10-bolditalic.otf]
\usepackage[margin=25mm]{geometry}
\usepackage[unicode,hidelinks]{hyperref}
\hypersetup{pdfauthor={LangevinLean Project}}
\setlength{\emergencystretch}{3em}
\setlength{\parskip}{0.4em}
\setlength{\parindent}{0pt}
\allowdisplaybreaks[2]
\raggedbottom
\title{\texorpdfstring{Local assumptions for differentiating relative entropy}{Local assumptions for differentiating relative entropy}}
\author{LangevinLean Project}
\date{}
\begin{document}\maketitle
Draft mathematical exposition; these notes may contain mistakes.


Let \(r_t\) be a family of positive densities relative to a probability measure \(\nu\). The local assumptions below give pointwise time derivatives and an integrable bound for \(\partial_t r_t(\log r_t+1)\) on a neighborhood of each time. They imply differentiability of the entropy integral. Normalization and logarithmic integrability identify that integral with relative entropy. An admissible weak equation then gives the entropy derivative estimate. The local assumptions, the weak equation and its energy conditions are retained as hypotheses.


\section{\texorpdfstring{The local differentiation assumptions}{The local differentiation assumptions}}


At a fixed time \(t\), assume a common time neighborhood on which the densities are positive and pointwise differentiable. Require measurability of the entropy and derivative integrands, integrability of the entropy integrand at \(t\), and an integrable spatial function bounding \(|\partial_u r_u(\log r_u+1)|\) throughout that neighborhood. These are local assumptions: the neighborhood and its bound may depend on \(t\). They contain no assumption about the derivative of the integrated entropy or its sign.


\section{\texorpdfstring{The derivative of the entropy integral}{The derivative of the entropy integral}}


Apply the scalar product and logarithm chain rules to \(r_u(x)\log r_u(x)\). The local bound in the hypotheses dominates this derivative uniformly for nearby \(u\). Dominated differentiation therefore proves integrability of \(\dot r_t(\log r_t+1)\) and gives

\[\frac d{dt}\int r_t\log r_t\,d\nu=\int\dot r_t(\log r_t+1)\,d\nu.\]

The conclusion follows from the stated density regularity and domination; differentiability of the entropy integral is not separately assumed.


\section{\texorpdfstring{An entropy derivative bound from the weak equation}{An entropy derivative bound from the weak equation}}


Write \(s_t=\nabla\log r_t\), \(I(t)=\int r_t\|s_t\|^2d\nu\), and \(E(t)=\int r_t\|e_t\|^2d\nu\). Local differentiation gives \(H'(t)=\int\dot r_t(\log r_t+1)d\nu\). Suppose the weak equation admits this test and both energies are finite. Its right-hand side is \(-I(t)+\int r_t\langle e_t,s_t\rangle d\nu\). Young's inequality bounds the last term by \(I(t)/4+E(t)\). Hence

\[H'(t)\le-\frac34I(t)+E(t).\]

This result uses the stated error energy but does not assert any numerical bound on that energy.


\section{\texorpdfstring{Finite relative entropy}{Finite relative entropy}}


At the selected time, the local assumptions give integrability of \(r_t\log r_t\). The positive \(C^1\) density is measurable, and normalization makes \(r_t\nu\) a probability measure. Changing measure gives integrability of \(\log r_t\) under \(r_t\nu\). Absolute continuity and the finite-divergence criterion then imply \(\operatorname{KL}(r_t\nu\mid\nu)<\infty\). This conclusion precedes the use of the real-valued entropy derivative.


\section{\texorpdfstring{Continuity of entropy on the stated interval}{Continuity of entropy on the stated interval}}


Suppose the local differentiation assumptions hold at every point of the closed interval. They imply local differentiability, and hence continuity, of \(t\mapsto\int r_t\log r_t\,d\nu\). At every time in the interval, positivity, normalization and integrability identify this integral with \(\operatorname{KL}(r_t\nu\mid\nu)\). Equality of the two functions on the interval transfers continuity to relative entropy, including continuity relative to the interval at its endpoints.


\section{\texorpdfstring{The right derivative of relative entropy}{The right derivative of relative entropy}}


Fix \(t\in[0,T)\). The local assumptions differentiate the ordinary entropy integral. For all sufficiently close \(u\ge t\), one has \(u\in[0,T]\); normalization and logarithmic integrability there identify the integral with relative entropy. Since these two functions agree at \(t\) and in a right neighborhood, they have the same right derivative at \(t\). This argument does not assume normalization or a probability interpretation outside the stated time interval.


\section{\texorpdfstring{The derivative estimate using relative Fisher information}{The derivative estimate using relative Fisher information}}


The change-of-density formula identifies \(\int\|\nabla\log r_t\|^2d(r_t\nu)\) with \(\int r_t\|\nabla\log r_t\|^2d\nu\). The corresponding integrability criterion proves that both are finite under the assumed score-energy bound. Thus the weighted energy in the weak-equation calculation is precisely the Fisher information used in LSI. Combining this identification with the right derivative of relative entropy gives \(H'(t)\le-3I(t)/4+E(t)\) for the same probability law.


\end{document}
